% !TEX program = xelatex
% The project report (issue #12). Every table is \input from _output/, every
% figure is \includegraphics from _output/, and every number quoted in prose
% is a macro from _output/report_values.tex (src/report_values.py). Nothing
% statistical is hand-typed. Compile: doit compile_latex_docs (requires both
% sample-period runs: `doit` and `SAMPLE_END=2024-12 doit`).

\newcommand*{\PathToOutput}{../_output}%

\documentclass[12pt]{article}
\usepackage{my_article_header}
\usepackage{my_common_header}
\usepackage{natbib}
\input{\PathToOutput/report_values.tex}

\title{The Virtue of Complexity in Return Prediction:\\ A Replication}
\author{Ahmed Alahmadi \and Jeff Key}
\date{\today}

\begin{document}
\maketitle

\begin{abstract}
\noindent We replicate the central out-of-sample results of
\citet{kelly2024virtue}: Figures~7, 8, and 11 at the 12-month training
window. Our replication matches the published anchor values at the ridgeless
high-complexity configuration within about one percent (alpha $t$-statistic
\MktAlphaT, information ratio \MktIR, Sharpe ratio \MktSharpe), reproduces
the paper's variable-importance ranking, and extends the analysis four ways:
updated data through 2024, government and corporate bonds, developed ex-US
equities, and the critique of \citet{nagel2025seemingly}. The virtue of
complexity replicates and survives updated data and bond markets; it does not
survive the trip abroad, and most of its performance is spanned by a simple
momentum rule.
\end{abstract}

\section{Introduction}

\citet{kelly2024virtue} make a provocative claim: in return prediction,
bigger models are better, without limit. They expand fifteen standard
predictors into up to $P = 12{,}000$ random Fourier features, train ridge
regressions on rolling windows of only $T=12$ months, and find that the
market-timing performance of the resulting strategy rises monotonically with
model complexity $c = P/T$, even though the out-of-sample $R^2$ of the
forecasts is negative everywhere. The claim matters because it contradicts
the intuition, standard since \citet{welch2008comprehensive}, that equity
premium prediction is too noisy to support rich models.

We replicate the paper's three headline exhibits at the 12-month training
window: Figure~7 (forecast accuracy and strategy moments against
complexity), Figure~8 (risk-adjusted trading performance against
complexity), and Figure~11 (variable importance). We then rerun everything
on data through 2024 and test external validity on bonds and international
equities, and we implement the momentum critique of
\citet{nagel2025seemingly}. Sections~\ref{sec:data}
through~\ref{sec:challenges} describe the data, the summary statistics, the
replication, the updated sample, the extensions, and the challenges.

\section{Data and conventions}\label{sec:data}

Predictors come from the monthly workbook maintained by Amit Goyal
\citep{welch2008comprehensive, goyal2024comprehensive}: fourteen of the
paper's fifteen predictors are constructed from it, and the fifteenth is the
lagged market excess return. The market return is the workbook's CRSP
value-weighted index return in excess of the T-bill rate; we cross-check it
against the CRSP index pulled from WRDS, and the two series agree almost
exactly. The bond extension uses the workbook's long-term government and
corporate bond returns; the international extension uses the developed
ex-US market factor from the Ken French data library.

We follow the paper's conventions exactly, because the replication turned
out to be sensitive to them. Returns are standardized by trailing 12-month
volatility and predictors by expanding-window volatility, both strictly
backward looking; the random-feature ridge is fit without an intercept; and
the random features are scaled within each training window. The authors
defend the volatility standardization and the no-intercept choice
explicitly in their reply \citep{kelly2025reply}. Two conventions deserve
emphasis because published sources do not pin them down and we selected
them by validating against published anchor values: the training-window
feature scaling is centered (standard deviations, not raw second moments),
and the market series is the workbook's, not the WRDS pull. The analysis
sample starts in 1930-01, after the standardization burn-in, and the paper
period ends in 2020-12, for \PaperMonths{} months. Relative to the paper we
economize on repetitions: 500 random-feature draws for the market grid
(the paper uses 1{,}000), 200 for the extensions, 100 for variable
importance, with each choice documented in the pipeline.

\section{Summary statistics}\label{sec:summary}

\begin{table}[htbp]\centering\footnotesize
\input{\PathToOutput/predictor_summary_table.tex}
\caption{The information set splits into slow and fast predictors. Panel~A
summarizes the raw market excess return and the fifteen predictors; Panel~B
the volatility-standardized versions the model consumes. The first-order
autocorrelations show the split that drives everything downstream: the
valuation ratios and yields are extremely persistent, while the return-type
predictors (the lagged market return, bond returns, and the default and
inflation series) move month to month. Standardization equalizes scales
without destroying that split.}
\label{tab:summary}
\end{table}

\begin{figure}[htbp]\centering
\includegraphics[width=\textwidth]{\PathToOutput/predictor_timeseries.png}
\caption{Standardized predictors at a glance. Slow series drift over
decades while fast series oscillate at monthly frequency; the model's
training window is twelve months, so the fast series are where a
twelve-month window can find variation to learn from. This is the visual
premise of the variable-importance result in Figure~\ref{fig:eleven}.}
\label{fig:timeseries}
\end{figure}

\section{Replication}\label{sec:replication}

\begin{figure}[htbp]\centering
\includegraphics[width=\textwidth]{\PathToOutput/figure7.png}
\caption{Complexity ruins the forecasts and it does not matter. Replication
of the paper's Figure~7: out-of-sample $R^2$ (Panel~A) is negative
everywhere and collapses at the interpolation boundary $c=1$, while the
coefficient norm, expected strategy return, and volatility (Panels~B
through~D) behave exactly as the paper's theory predicts. Panels in return
units sit about a quarter below the paper's levels, a constant
forecast-scale difference documented in the pipeline; every scale-free
statistic matches.}
\label{fig:seven}
\end{figure}

\begin{figure}[htbp]\centering
\includegraphics[width=\textwidth]{\PathToOutput/figure8.png}
\caption{The virtue of complexity, replicated. The paper's crux figure:
Sharpe ratio, alpha, information ratio, and alpha $t$-statistic all rise
with complexity for every shrinkage level, despite the negative $R^2$ in
Figure~\ref{fig:seven}. At the ridgeless $c=1{,}000$ anchor our replication
delivers a Sharpe ratio of \MktSharpe, an information ratio of \MktIR, and
an alpha $t$-statistic of \MktAlphaT, each within about one percent of the
published anchor values.}
\label{fig:eight}
\end{figure}

The contrast between Figures~\ref{fig:seven} and~\ref{fig:eight} is the
finding, not a bug: forecast accuracy measured by $R^2$ is useless while
the trading strategy built from the same forecasts is strong, because the
strategy rewards getting the sign and relative magnitude right, not the
level.

\begin{figure}[htbp]\centering
\includegraphics[width=0.9\textwidth]{\PathToOutput/figure11.png}
\caption{Twelve months of data is enough because the model feeds on fast
variables. Replication of the paper's Figure~11: variable importance is the
loss in out-of-sample $R^2$ (bars) and Sharpe ratio (line) from removing
each predictor. The lagged market return dominates (removing it costs
\VILeaderRTwoPct{} percentage points of $R^2$), followed by the long-term
bond return, the default return, and inflation, all fast-moving series;
the persistent valuation ratios contribute nothing detectable, matching
the paper's ranking.}
\label{fig:eleven}
\end{figure}

\section{Updated sample}\label{sec:updated}

Every artifact in the project regenerates on data through 2024-12
(\UpdMonths{} months) from a single configuration change. The finding
survives: at the ridgeless anchor the Sharpe ratio is \MktSharpeUpd{}
(paper period \MktSharpe), the information ratio \MktIRUpd{} (\MktIR), and
the alpha $t$-statistic \MktAlphaTUpd{} (\MktAlphaT). Performance softens
slightly with four more years of data but the rising-with-complexity shape
is unchanged, as Figure~\ref{fig:eightupd} shows.

\begin{figure}[htbp]\centering
\includegraphics[width=\textwidth]{\PathToOutput/figure8_2024-12.png}
\caption{The virtue of complexity survives updated data. The Figure~8
panels re-estimated on 1930 through 2024: every curve retains its shape,
with levels slightly below the paper period, so the paper's result is not
an artifact of its sample end.}
\label{fig:eightupd}
\end{figure}

\section{Extensions}\label{sec:extensions}

\subsection{Bonds}

We point the identical machinery at long-term government and corporate bond
excess returns, both built from workbook columns, with each target's own
lag replacing the market's.

\begin{figure}[htbp]\centering
\includegraphics[width=\textwidth]{\PathToOutput/figure_bonds.png}
\caption{The virtue of complexity generalizes to bonds. For both bond
markets the out-of-sample $R^2$ stays useless while the Sharpe ratio rises
with complexity. Corporate bonds (ridgeless Sharpe \BondCorpSharpe) come
close to the equity market (\MktSharpe), consistent with the information
set's rich credit content; pure duration risk (government bonds, ridgeless
Sharpe \BondGovSharpe) supports less. Both conclusions hold on updated data
(\BondCorpSharpeUpd{} and \BondGovSharpeUpd).}
\label{fig:bonds}
\end{figure}

\subsection{International equities}

As an external-validity test we forecast the developed ex-US market excess
return from the Ken French library with the same US information set. Two
caveats are built into the design: US predictors are being asked to
forecast non-US returns, and the series only starts in 1990, leaving about
a third of the market study's months.

\begin{figure}[htbp]\centering
\includegraphics[width=\textwidth]{\PathToOutput/figure_intl.png}
\caption{The shape generalizes abroad, the level does not. Sharpe ratios
still rise with complexity, but they top out at \IntlSharpeBest{} for the
best shrinkage level (ridgeless \IntlSharpe), a fraction of the US levels,
and statistically indistinguishable from zero on the short sample. On
updated data the ridgeless level falls further, to \IntlSharpeUpd. With US
predictors and three decades of data there is no measurable international
timing skill, only the qualitative complexity pattern.}
\label{fig:intl}
\end{figure}

% Nagel critique subsection (issue #17), \input by report_kmz.tex.
% Generated artifacts consumed here:
%   \input{\PathToOutput/nagel_comparison_table.tex}   (from src/table_nagel.py)
%   \PathToOutput/figure_nagel.png
% No hand-typed statistics: prose numbers come from report_values.tex macros.

\subsection{Is the virtue of complexity just momentum? Nagel's critique}

\citet{nagel2025seemingly} argues that when the number of random features $P$
greatly exceeds the training window $T$, the ridgeless RFF forecast collapses
to a Gaussian-kernel regression: a similarity-weighted average of the $T=12$
training-window returns. Because the standardized predictors are persistent,
similarity is approximately recency, and kernel distances scale with predictor
volatility, so the strategy is, mechanically, recency-weighted and
volatility-timed momentum. We run his three tests on our own pipeline at the
anchor configuration (the ridgeless model at complexity $c=1{,}000$) and
report what we find, whichever side it favors.

Table~\ref{tab:nagel} scores our replicated VoC strategy against Nagel's
hand-built momentum benchmark, which places linearly declining weights on the
trailing twelve standardized returns, on identical out-of-sample months and
identical metrics (the same performance module used everywhere else in the
project). The benchmark's Sharpe ratio of \NagelBenchSharpe{} sits next to the
VoC row's \NagelVoCSharpe: a strategy with no machine learning in it earns
almost the same risk-adjusted return. One convention note: the VoC row scores
the across-seed mean forecast path, which denoises the forecast slightly, so
it sits a touch above the average-the-statistics value (\MktSharpe) that the
replication reports elsewhere; the comparison within the table is
apples-to-apples.

\begin{table}[htbp]\centering
\input{\PathToOutput/nagel_comparison_table.tex}
\caption{The virtue of complexity is largely spanned by simple momentum. The
VoC strategy and Nagel's recency-weighted momentum benchmark, scored on the
same out-of-sample months with the same metric definitions, deliver similar
performance, and the spanning regression of the VoC strategy on the benchmark
and the static market leaves an alpha with a $t$-statistic of
\NagelSpanAlphaT, below conventional significance.}
\label{tab:nagel}
\end{table}

Figure~\ref{fig:nagel} anatomizes the VoC forecast by regressing it on the
trailing twelve standardized returns; the estimated lag weights are shown
against the benchmark's declining profile. The regression explains only
$R^2 = \NagelAnatomyRTwo$ of the forecast's variation, so the forecast
contains more than recency-weighted momentum, even though its tradable
performance is largely spanned by it.

\begin{figure}[htbp]\centering
\includegraphics[width=0.8\textwidth]{\PathToOutput/figure_nagel.png}
\caption{The VoC forecast is momentum-tilted but not momentum. Estimated lag
weights of the VoC forecast on the trailing twelve standardized returns
(regression $R^2 = \NagelAnatomyRTwo$) against the benchmark's linearly
declining weights: the profiles share the recency tilt, but most of the
forecast's variation is left unexplained.}
\label{fig:nagel}
\end{figure}

\paragraph{The authors' reply.} \citet[Section~4.1]{kelly2025reply} respond
that the momentum reading is a feature, not an artifact: the volatility
standardization is economically motivated, and a reduction toward a
recency-weighted average of past returns is a property any well-specified
nonparametric return predictor would share, not evidence that the
out-of-sample gains are spurious. We present both sides and let the
replicated numbers speak: the performance is real but a much simpler
strategy captures most of it.


\section{Challenges and lessons}\label{sec:challenges}

Three challenges shaped the project. First, compute: the naive primal ridge
at $P = 12{,}000$ is infeasible on a laptop, so the engine works in the dual
with one eigendecomposition per training window shared across all shrinkage
levels, which makes the full 500-seed grid a minutes-scale job. Second,
conventions: our first implementation matched the paper's shapes but missed
the published anchor values by a wide margin, and closing the gap required a
systematic investigation that identified the centered feature scaling and
the workbook market series as the paper's conventions; the anchors now serve
as permanent regression tests. Third, honesty about residuals: strategy
moments in return units sit at about three quarters of the paper's levels, a
constant forecast-scale difference that cancels in every ratio statistic; we
document it rather than tune it away, since no published scale-free anchor
pins it down. The one result we could not and did not try to rescue is the
international study, where the honest conclusion is a null.

\bibliographystyle{jpe}
\bibliography{bibliography}

\end{document}
